\section{The formal model and its accepted language}
We now state the object precisely enough that a theorem has a definite scope.
Let $\Sigma$ be a finite alphabet, $\Sigma^*$ its finite words and
$\varepsilon$ the empty word. Positions in code are zero-based; the letter
$u_r$ in mathematics is the $r$th letter, with $1\leq r\leq |u|$.

\begin{definition}[Paired domain and advancing automaton]
Let $\rho\subseteq\Sigma\times\Sigma$ be symmetric. Define
\[
 \operatorname{WK}_{\rho}=
 \{(u,v):|u|=|v|,\ (u_r,v_r)\in\rho\text{ for every }r\}.
\]
An advancing Watson--Crick automaton is
$M=(\Sigma,\rho,Q,q_0,F,\mathcal R)$, where $Q$ and $\mathcal R$ are finite,
$q_0\in Q$, $F\subseteq Q$, and
$\mathcal R\subseteq Q\times\Sigma^*\times\Sigma^*\times Q$.
Each rule $(q,x,y,q')$ satisfies $|x|+|y|>0$.
\end{definition}
The word \emph{advancing} names our restriction; general literature permits
additional conventions. This presentation uses the paired-input framework of
\cite{czeizler}, but proves the following constructions under the definitions
given here, rather than importing a language hierarchy wholesale.

For fixed $(u,v)\in\operatorname{WK}_\rho$ of length $n$, write
\begin{equation}
 (q,i,j)\vdash_{u,v}(q',i+|x|,j+|y|)
 \label{eq:stepformal}
\end{equation}
when $(q,x,y,q')\in\mathcal R$, $x$ is a prefix of $u[i:]$ and $y$ is a
prefix of $v[j:]$. The bounds on cursors follow from prefix matching.
The reflexive transitive closure $\vdash^*$ permits zero steps.
\begin{definition}[Pair acceptance and upper language]
The accepted relation and its upper projection are
\begin{align}
 A(M)&=\{(u,v)\in\operatorname{WK}_\rho:
       (q_0,0,0)\vdash_{u,v}^*(f,n,n),\ f\in F\},\\
 L(M)&=\{u:\exists v\ (u,v)\in A(M)\}.
\end{align}
\end{definition}
There are two existential choices: a lower word and a legal run on that
fixed pair. A program that decides pair acceptance does not automatically
decide the upper language without resolving the first choice. The empty
word belongs to $L(M)$ exactly when $q_0\in F$, because advancing rules
cannot consume anything on empty input.

The relation constrains aligned \emph{input} positions, not arbitrary blocks
consumed at different cursor positions. If $i\ne j$, a rule can consume
upper a and lower b even when $\rho$ is identity. Requiring the two rule
labels themselves to form a paired word would incorrectly forbid the
equal-block construction. This is why the transition notation must not be
mistaken for a chemical hybridization instruction.

\subsection{Synchronous reading recovers ordinary regularity}
\begin{proposition}[Synchronous single-letter restriction]
If every rule consumes exactly one letter on each strand, then $L(M)$ is
regular. Conversely, every regular language has such an automaton with
identity pairing.
\end{proposition}
\begin{proof}
Construct a nondeterministic finite automaton on upper letters, with the same
states, start state and final states. Include an edge $q\xrightarrow{a}q'$
exactly when some $b$ satisfies $(a,b)\in\rho$ and
$(q,a,b,q')\in\mathcal R$. Each paired run projects to an NFA run.
Conversely, choose one witnessing $b$ for each edge of an NFA run. Their
concatenation forms an admissible lower word and a paired run, because both
cursors advance together. The accepted languages coincide. For the reverse
direction, replace every edge $q\xrightarrow{a}q'$ of an NFA by
$(q,a,a,q')$ and use identity pairing.
\end{proof}
Thus double-stranded notation alone does not explain nonregular recognition.
The distinguishing resource in our equal-block example is asynchronous access.
This proposition concerns single-letter synchronous rules; it does not assert
that every variant of a multihead automaton has the same characterization.

\section{Three levels of determinism}
\index{Weak determinism}\index{Syntactic determinism}\index{Strong determinism}
Write $x\preceq z$ when $x$ is a prefix of $z$, and $x\sim_p y$ when
one is a prefix of the other. In particular, $\varepsilon$ is comparable
with every word. The terminology below follows \cite{czeizler}.

\begin{definition}[Determinism conventions]
\emph{Weak determinism} means at most one enabled rule in every reachable
configuration on admissible paired inputs. \emph{Syntactic determinism}
requires, for distinct rules $(q,x,y,q')$ and $(q,x',y',q'')$,
\begin{equation}
 x\not\sim_p x'\quad\text{or}\quad y\not\sim_p y'.
 \label{eq:det}
\end{equation}
\emph{Strong determinism} adds identity pairing to syntactic determinism.
\end{definition}
The word syntactic distinguishes the check on written rules from the semantic
property of reachable executions. A unique run for a fixed pair also does not
imply a unique lower word for an upper input under a nonfunctional $\rho$.

\begin{proposition}[A sufficient local test]
Equation~\ref{eq:det} implies weak determinism.
\end{proposition}
\begin{proof}
If two rules were enabled, their upper labels would both prefix the same
unread upper suffix, hence would be prefix-comparable. The same argument
would apply to the lower labels. This contradicts Equation~\ref{eq:det}.
No assumption about how that configuration was reached is needed.
\end{proof}

\Plate{06-determinism}{Two outgoing rules conflict syntactically when both
label pairs are prefix-comparable. The unread pair shown enables both rules.
In general a reported conflict need not be reachable under a particular
pairing relation; this diagram demonstrates one concrete admissible case.}{fig:determinism}

\Listing{prefix_conflicts}
For the equal-block rules, the conflict list is empty. In seek, the upper
labels $\varepsilon$ and a are comparable, but lower a and b are not.
In pair, upper a and b are not comparable. Drain has one rule. Together
with identity pairing, this proves that the construction is strongly
deterministic under our conventions. Calling its implementation a search
does not make the recognized computation nondeterministic.

Failure of the test does not establish reachable ambiguity. For example,
append an unreachable state with rules $(z,a,a,z)$ and $(z,aa,aa,z)$.
They conflict syntactically, yet no reachable run is changed. This is a
counterexample to necessity, not a general algorithm for deciding weak
determinism. The checker reports pairs of rule indices; it deliberately
makes no claim to solve that semantic decision problem.

\section{Normal forms and honest descriptional complexity}
\index{Normal form}\index{Descriptional complexity}
\begin{proposition}[One-symbol normalization]
Every advancing automaton has an equivalent advancing automaton whose rules
consume one symbol on exactly one strand. Equivalence here means equality
of accepted paired relations, with the same empty-input convention.
\end{proposition}
\begin{proof}
For a rule consuming $(x,y)$, let $\ell=|x|+|y|>0$. Replace it by a private
path of $\ell$ transitions, consuming the letters of $x$ on the upper
strand first and then the letters of $y$ on the lower. Introduce $\ell-1$
fresh nonfinal states used only by this path. Keep original start and final
states. A successful old step expands to a successful path. Conversely,
an accepting new run cannot end inside a private path and cannot switch
between private paths, so contract each complete path to its original rule.
Both constructions preserve the consumed words and endpoints.
\end{proof}
Matching the upper block and then failing to match the lower block creates
a dead branch, not a new accepting run. Private states are essential: sharing
them indiscriminately could join the first half of one rule to the second
half of another.
\Listing{expand_rules}
Pass all declared states as \texttt{reserved} when they include isolated
start or final states absent from the rule endpoints. The generated private
names must be fresh relative to the entire state set, not just reachable
rule sources. The code validates advancing rules before expanding them.

This elementary construction preserves nondeterministic acceptance, not
syntactic determinism. Consider $(s,a,b,f)$ and $(s,a,a,g)$: lower b and
a distinguish the original rules. Expanding upper letters first gives two
outgoing $(a,\varepsilon)$ rules. A theorem about deterministic normal forms
requires a more careful construction; it cannot be inferred from this code.

For rule lengths $\ell_r$, this construction has
\begin{equation}
 |Q'|=|Q|+\sum_{r\in\mathcal R}(\ell_r-1),\qquad
 |\mathcal R'|=\sum_{r\in\mathcal R}\ell_r.
\end{equation}
Report at least $(|Q|,|\mathcal R|,\sum_r\ell_r,|\rho|)$ when comparing
descriptions. A rule storing a thousand symbols is not comparable to a
single-symbol rule merely because both count as one arrow. Normalization
trades label storage for explicit control states; it does not erase cost.
